\documentclass[11pt,a4paper]{article}
\usepackage[a4paper,margin=2.4cm]{geometry}
\usepackage{amsmath,amssymb}
\usepackage{xcolor}
\usepackage{enumitem}
\usepackage{parskip}
\usepackage[hidelinks]{hyperref}

\title{Response to the reviewers\\
\large Manuscript NEXRES-D-26-02183\\
\textit{A Four-Anchor Phenomenological Fit of a $1/\ln^2(t/t_{\mathrm{P}})$ Hubble Relaxation}}
\author{Liang Wang}
\date{Draft for the revision due 27 November 2026. Not yet resubmitted.}

\definecolor{revcolor}{HTML}{444444}
\newenvironment{review}{\par\color{revcolor}\itshape\leftskip=1.2em\rightskip=0.4em\noindent}{\par\leftskip=0pt\rightskip=0pt}
\newcommand{\revtodo}[1]{\textcolor{red}{\textbf{TODO:} #1}}

\begin{document}
\maketitle

Dear Dr Biswas,

Thank you for handling manuscript NEXRES-D-26-02183 and for the reports from both reviewers.
The revision follows Reviewer~1 and narrows the paper.
The point-by-point replies are below.
Locations refer to sections of the revised manuscript.
Page numbers are omitted because the draft has not been frozen for production.
Items that need a refit are marked in red.
No new fitted value has been invented to fill those items.

\section*{Summary of the revision}

\begin{itemize}[leftmargin=1.4em]
  \item The core empirical result is the four-anchor Hubble fit (Pillar~C), kept with its quoted summary $H_\infty = 104.5\,\mathrm{km\,s^{-1}\,Mpc^{-1}}$, $\beta = -6.24\,\mathrm{km\,s^{-1}\,Mpc^{-1}}$, $R^2 = 0.91$, and $\chi^2/\mathrm{dof} = 0.28$ on two degrees of freedom.
  The working equation is $H(t) = H_\infty + \beta\,10^{5}/\ln^2(t/t_{\mathrm{P}})$, which is the normalization in the figure and forecast notebooks.
  That factor reconciles the intercept $104.5$ with the quoted finite-redshift values near $71$--$72$.
  \item The Riemann-zeta motivation has been removed.
  The fitting function is described as a selected phenomenological ansatz.
  No uniqueness claim remains.
  \item The lattice CMB forward model, the inverse reconstruction, and the acoustic freeze-out discussion have been removed from the main text and from the Supplementary Material (\texttt{sections/S\_cmb\_lattice.tex} is withdrawn and is not input).
  There is no three-scale validation claim.
  \item The Keck, VLT, and full-sample quasar comparison is exploratory.
  Detection language has been removed.
  A side-by-side table of the submitted quotations is in Section~4.2.
  The clash between $\Gamma = +6.42\pm 1.14$ and $\Gamma = +5.56\pm 0.99$ is left unresolved.
  \item Pillar~M now states a falsification standard.
  A $\chi^2/\mathrm{dof}$ of order $10^{3}$ on DESI DR2 is described as a poor fit.
  \item The Elsevier wrapper journal line is Next Research.
  The title and abstract have been rewritten for the narrowed scope.
  \item \revtodo{Still required before resubmission: one Murphy-sample refit; the anchor-list confirmation, covariance, and leave-one-anchor-out table; a common-likelihood AIC/BIC comparison of the ansatz, flat $\Lambda$CDM, and CPL; a forecast with a stated covariance; pinned software versions and one constant table.}
\end{itemize}

\section*{Reviewer 1}

\subsection*{R1.1}
\begin{review}
Internal inconsistency in the headline result ($\Gamma = +6.42\pm 1.14$ in Section~4.2 versus $+5.56\pm 0.99$ in the abstract, the conclusion, and Table~1).
\end{review}

The two quotations are both still visible, and neither has been chosen.
Section~4.2 of the submitted manuscript quotes $\Gamma = +6.42\pm 1.14$ and $\Delta\chi^2 = 31.54$ for the Murphy sample under the fit-error budget.
The instrument table in that section, the submitted abstract, and one conclusion sentence quote $\Gamma = +5.56\pm 0.99$.
The conclusion also quotes $\Delta\chi^2 = 31.83$.
The revised Section~4.2 prints this clash and withdraws every detection claim that depended on either number.
The quasar material is labeled exploratory.

\revtodo{Refit the Murphy et al.\ (2003) sample once, under one stated error budget, and replace both quotations with that single result. Do not adopt $+6.42\pm 1.14$ or $+5.56\pm 0.99$ by hand.}

Change: Abstract; Section~4.2 and Table of quoted subset coefficients; Conclusion.

\subsection*{R1.2}
\begin{review}
Subset dependence of the quasar drift signal (Keck-only, null on the full King 2012 catalog, sign reversal in the VLT-only subset).
\end{review}

Section~4.2 and the Supplementary Material place the full King et al.\ (2012) catalog, the Keck-only subset, and the VLT-only subset side by side, using the coefficients quoted in the submitted manuscript.
On those quotations the Keck-only sign is positive, the VLT-only sign is negative, and the full catalog was quoted with $\Delta\mathrm{BIC} = +5.7$ in favor of a constant.
The text states that no pre-registered subset protocol exists, and it leaves both readings open: an instrumental systematic, and a subsample fluctuation.
Words such as detection, confirmation, and cross-validation have been removed from this comparison.

\revtodo{After the R1.1 refit, recompute the three columns (full sample, Keck-only, VLT-only) under the same budget. The present table is a record of the submitted quotations.}

Change: Section~4.2; Supplementary figures for the King catalog and the subset scan.

\subsection*{R1.3}
\begin{review}
Non-falsifiable framing of the leverage argument (Pillar~M); $\chi^2/\mathrm{dof}\sim 10^{3}$ conflated with resolvability.
\end{review}

Section~4.4 keeps the arithmetic range of $10^{5}/\ln^2(t/t_{\mathrm{P}})$ as a property of the selected function and separates it from goodness of fit.
The falsification standard is stated there: the four-anchor description fails if a leave-one-anchor-out refit misses the omitted anchor, or if a direct $H(z)$ measurement with a stated covariance falls outside the curve; each dense sample fails separately when its own $\chi^2$ is unacceptable.
A non-detection where the function barely moves is recorded as a non-detection.
A detection in that window would be a failure of the expectation that the function is unresolved there.
Section~4.3 calls the DESI DR2 value $\chi^2/\mathrm{dof}\sim 10^{3}$ a poor fit.
The conclusion does not list that value as support.

\revtodo{The leave-one-anchor-out half of the standard has no numbers yet. Insert them only after the refit in R2.5.}

Change: Sections~4.3 and~4.4; Discussion; Conclusion.

\subsection*{R1.4}
\begin{review}
Retraction of the lattice CMB forward model and inverse reconstruction, and the resulting gap in the three-scale claim.
\end{review}

This material has been removed from the revision, as Reviewer~1 recommended.
The main text and the Supplementary Material no longer contain the lattice forward model, the inverse reconstruction, the held-out hemisphere test, or the acoustic freeze-out discussion.
\texttt{sections/S\_cmb\_lattice.tex} is a withdrawal note and is not input.
The title, abstract, and conclusion contain no three-scale validation claim.
What remains as the core empirical result is the four-anchor fit, with the limitations in Section~6.

Change: entire manuscript; Supplementary Material; title and abstract.

\section*{Reviewer 2}

\subsection*{R2.1}
\begin{review}
Insufficient physical justification of the proposed cooling law.
\end{review}

Section~2 states that $1/\ln^2(t/t_{\mathrm{P}})$ is a selected phenomenological ansatz.
The slow variation of the logarithm is the reason this shape was tried.
The Riemann-zeta argument has been removed, and no microscopic derivation is claimed.
Applicability is limited to the comparisons actually performed: a curve through four assigned Hubble anchors, an exploratory quasar parametrization, and the late-time $\chi^2$ values reported as fits.

Change: Sections~1 and~2.

\subsection*{R2.2}
\begin{review}
The theoretical derivation in Section~2.2 requires further mathematical justification.
\end{review}

The spectral-density derivation has been deleted rather than repaired.
Section~2 defines $\xi(t) = 1/\ln^2(t/t_{\mathrm{P}})$, writes the normalized Hubble equation used by the notebooks, and differentiates the fine-structure parametrization.
The dimensions of $H_\infty$, $\beta$, and the cosmological term are discussed in Section~2.1.
No step in the revised text identifies a vacuum-energy variance with the square of a level spacing.

Change: Section~2. The former Section~2.2 is gone.

\subsection*{R2.3}
\begin{review}
Uniqueness of the logarithmic cooling law is insufficiently supported.
\end{review}

Uniqueness language has been removed from the abstract, the introduction, the ansatz section, and the conclusion.
Section~2.3 says explicitly that other slow functions, including $1/\ln^p(t/t_{\mathrm{P}})$ for $p\neq 2$, are not excluded.
No uniqueness theorem is offered.

Change: Abstract; Sections~1, 2, and~7.

\subsection*{R2.4}
\begin{review}
Inconsistency in the physical interpretation of the Hubble parameter, including $H_\infty = 104.5\,\mathrm{km\,s^{-1}\,Mpc^{-1}}$ and the high-redshift behavior.
\end{review}

The revised equation is $H(t) = H_\infty + \beta\,10^{5}/\ln^2(t/t_{\mathrm{P}})$.
The factor $10^{5}$ is the normalization in \texttt{notebooks/fig19\_jwst\_h0\_anchors.ipynb} and in \texttt{notebooks/fig31\_pillar\_O\_forecast.ipynb}.
With the quoted pair $(H_\infty,\beta) = (104.5,-6.24)$ and the submitted leverage value $10^{5}\xi\approx 5.19$ at $z\approx 2$, the same arithmetic gives $H\approx 72.1\,\mathrm{km\,s^{-1}\,Mpc^{-1}}$, which is the quoted $H(z=2)$.
So $H_\infty$ is the intercept at $t\to\infty$ inside this fitting function, and the values near $71$--$72$ are finite-$t$ evaluations of that function.
Because $\beta<0$, the curve increases toward $H_\infty$.
The submitted claim $H_\infty < H_0$ does not describe the fit.
The identification $H_\infty = \sqrt{\Lambda_\infty/3}$ is withdrawn: $H_\infty$ is a fit coefficient in $\mathrm{km\,s^{-1}\,Mpc^{-1}}$, and that equality is not dimensionally imposed.

A separate inconsistency remains.
The submitted prose named the anchors as SH0ES, JWST, Planck, and BAO.
The figure notebook assigns Planck $67.4\pm 0.5$ at $3.8\times 10^{5}\,\mathrm{yr}$, a TRGB point $70.4\pm 1.9$ at $1.5\times 10^{9}\,\mathrm{yr}$, a Mira point $72.37\pm 2.97$ at $6.0\times 10^{9}\,\mathrm{yr}$, and SH0ES $73.04\pm 1.04$ at $1.38\times 10^{10}\,\mathrm{yr}$.
Several of these are present-day $H_0$ inferences placed at earlier times.
The Planck number is a $\Lambda$CDM inference of today's Hubble constant, not the expansion rate at recombination.
Section~4.1 states that the fit is a curve through those assigned pairs.

\revtodo{Confirm which four anchors are the sample, attach a checked citation to the TRGB and Mira rows, and recompute $H(z)$ from one $t(z)$ only if the confirmed parameters change the quoted $71$--$72$ sequence. Also choose one convention for $(\Lambda_\infty,\gamma)$.}

Change: Sections~2.1, 3.1, 4.1, and~5.1.

\subsection*{R2.5}
\begin{review}
Statistical reliability of the four-anchor fit: degrees of freedom, covariance, and leave-one-anchor-out.
\end{review}

Section~3.1 records the definitions used in the notebook.
$R^2$ is unweighted.
$\chi^2$ uses the quoted errors.
There are four points and two parameters, so $\mathrm{dof} = 2$.
The notebook calls \texttt{curve\_fit} with \texttt{absolute\_sigma=True} and stores a covariance that the manuscript never printed.
No leave-one-anchor-out loop is in the notebook.
Section~4.1 states that an $R^2$ of $0.91$ and a reduced $\chi^2$ of $0.28$ on two degrees of freedom do not establish an expansion history.

\revtodo{Print the $2\times 2$ covariance and the four leave-one-anchor-out refits after the anchor list is confirmed. Do not insert a covariance by hand.}

Change: Sections~3.1 and~4.1.

\subsection*{R2.6}
\begin{review}
A poor goodness of fit must not be treated as support.
\end{review}

This is the same change as R1.3.
Section~4.3 and Table of late-time $\chi^2$ values present the quoted $\chi^2/\mathrm{dof}$ as goodness of fit.
The DESI DR2 entry of order $10^{3}$, the joint entry $23.7$, the cosmic-chronometer entry $4.01$, and the Pantheon+ entry $1.27$ are worse descriptions than the quoted flat $\Lambda$CDM values.
They are not counted as evidence for the ansatz.
The conclusion separates that fact from the four-anchor number.

Change: Sections~4.3, 4.4, and~7.

\subsection*{R2.7}
\begin{review}
Inadequate comparison with $\Lambda$CDM and other models on the same datasets.
\end{review}

The revised text reports only the $\chi^2$ values already printed in the submitted manuscript, and it labels them as previously quoted.
It does not treat those values as a finished model comparison.
The submitted $\Delta\mathrm{BIC} = -2.6$ between CPL and flat $\Lambda$CDM is not reused as a comparison against the ansatz.

\revtodo{On one stated dataset and covariance, report the likelihood, $\chi^2$, AIC, and BIC for the ansatz, flat $\Lambda$CDM, and CPL, with the parameter count of each. Do not invent those information criteria in the response.}

Change: Sections~3.3 and~4.3; Supplementary late-time captions. The numerical comparison itself is still a to-do.

\subsection*{R2.8}
\begin{review}
Uncertainty analysis of the fine-structure variation is insufficient.
\end{review}

Section~4.2 keeps both submitted error budgets side by side and states that the formal ratio changes between them.
The text no longer promotes either budget to a detection.
Instrumental calibration and absorber-to-absorber correlations beyond the two budgets already quoted are not a new error budget.

\revtodo{Write one error budget that states fit error, random term, calibration systematic, and any covariance, and refit under that budget. The present table only repeats the submitted budgets.}

Change: Sections~3.2 and~4.2.

\subsection*{R2.9}
\begin{review}
Opposite signs from Keck-only and VLT-only observations.
\end{review}

This is answered with R1.2.
The revised side-by-side table shows the opposite signs on the submitted quotations.
The discussion cites the wavelength-calibration papers as one available reading and states that a subset fluctuation is also still available.
The combined rows are closer to zero than the Keck-only union on those quotations.
No cosmological interpretation is adopted.

Change: Section~4.2; Discussion; Supplementary subset figure.

\subsection*{R2.10}
\begin{review}
Atomic-clock constraints need a derivation of the present-day rate and its parameter dependence.
\end{review}

Section~2.2 differentiates $\Delta\alpha/\alpha = \Gamma\times 10^{-5}\,[\xi(t)-\xi(t_{\mathrm{now}})]$ to
\[
  \dot\alpha/\alpha = \Gamma\times 10^{-5}\times(-2)/\big(t\ln^3(t/t_{\mathrm{P}})\big).
\]
The submitted prefactor $5.25\times 10^{-22}\,\mathrm{yr}^{-1}$ per unit $\Gamma$ is the present-epoch evaluation of that derivative for $\ln(t_{\mathrm{now}}/t_{\mathrm{P}})\approx 140$.
Dividing the Rosenband et al.\ bound $1\times 10^{-17}\,\mathrm{yr}^{-1}$ by that prefactor gives the conditional limit $|\Gamma|\lesssim 1.9\times 10^{4}$.
The limit depends on $\Gamma$ only as the amplitude of this parametrization.
It does not depend on $\beta$, and it does not use the unreconciled quasar value of $\Gamma$.
No joint posterior of the laboratory bound and the quasar catalogs is constructed.
Section~4.5 states that dependence.

\revtodo{If the R1.1 refit changes the quasar $\Gamma$, update only the illustration that overlays those points on the laboratory translation. The translation itself does not need a new $\Gamma$.}

Change: Sections~2.2 and~4.5.

\subsection*{R2.11}
\begin{review}
Cosmological stability analysis is incomplete.
\end{review}

The five background checks remain only as a previously quoted scan, in the Supplementary Material, with a caption that a pass is not a stability proof.
Perturbative stability, gradient instabilities, and the origin of the prescribed sound speed are listed as not done (Section~6).
The main text does not claim that the ansatz is stable.

\revtodo{A perturbative and gradient analysis is outside this narrowed revision. If it is added later, it needs its own calculation and should not be filled in from the five background checks.}

Change: Section~4.3; Supplementary stability caption; Section~6.

\subsection*{R2.12}
\begin{review}
The connection between lattice dynamics and cosmological observables is unclear.
\end{review}

This material has been removed from the revision.
The lattice field is not used as a cosmological observable anywhere in the revised main text or Supplementary Material.

Change: removed. Former lattice sections and \texttt{sections/S\_cmb\_lattice.tex}.

\subsection*{R2.13}
\begin{review}
Symplecticity of the damped numerical scheme requires clarification.
\end{review}

This material has been removed from the revision.
The revised manuscript does not describe a St\"ormer--Verlet scheme, a drag term, or a symplectic map.

Change: removed.

\subsection*{R2.14}
\begin{review}
Numerical convergence and grid-independence tests are necessary.
\end{review}

This material has been removed from the revision.
There is no lattice resolution, time step, or grid in the revised manuscript.
Convergence of the statistical fits is a separate item, marked under R2.19.

Change: removed.

\subsection*{R2.15}
\begin{review}
Discrepancy between simulated and observed CMB acoustic peaks.
\end{review}

This material has been removed from the revision.
No simulated $C_\ell$, acoustic-peak table, or Planck-spectrum $\chi^2$ is reported, and none is cited as a limitation of a model that is no longer in the paper.

Change: removed.

\subsection*{R2.16}
\begin{review}
The inverse reconstruction procedure requires stronger validation.
\end{review}

This material has been removed from the revision.
The in-sample correlation and the held-out hemisphere test are not cited.

Change: removed.

\subsection*{R2.17}
\begin{review}
The treatment of acoustic freeze-out requires further explanation.
\end{review}

This material has been removed from the revision.
No freeze-out scale, threshold, or sound-horizon analog is defined.

Change: removed.

\subsection*{R2.18}
\begin{review}
Forecasted observational discrimination appears overly optimistic.
\end{review}

Section~3.3 documents the existing mock from the forecast notebook: a flat $\Lambda$CDM age integral with $H_0 = 67.4$ and $\Omega_m = 0.315$, independent Gaussian fractional errors, 30 realizations, Generator seed 42, and Nelder--Mead refits.
The quoted figure of merit is $\sqrt{\Delta\chi^2}$ of the median draw.
Section~4.6 copies the submitted table and states that these entries are not survey significances.
The abstract and conclusion do not claim an $11\sigma$ forecast.
The separation between the fitting function and flat $\Lambda$CDM at $z>2$ is described as a property of the two functions, and the ability of a real catalog to resolve it is left to a covariance that the mock does not have.

\revtodo{Replace the table with a forecast that states the redshift sampling, the covariance, and the systematic floor. Do not keep $\gtrsim 11\sigma$ as a predicted significance.}

Change: Sections~3.3, 4.6, and~7.

\subsection*{R2.19}
\begin{review}
Reproducibility of the numerical and statistical analysis.
\end{review}

Section~3.4 names the repository and the notebooks that generate the retained figures.
It records Python~3.12.3 from the Pillar~C notebook metadata, \texttt{scipy.optimize.curve\_fit} with \texttt{absolute\_sigma=True} for the anchor fit, and the forecast seed, realization count, and optimizer settings above.
It also records that the anchor notebook and the forecast notebook do not share one value of $t_{\mathrm{P}}$ or one seconds-per-year conversion.

\revtodo{Pin NumPy and SciPy versions, adopt one constant table, and state bounds and convergence tolerances for every quoted fit. Do not guess the unpinned versions.}

Change: Section~3.4.

\subsection*{R2.20}
\begin{review}
Statistical model selection and the conversion of likelihood ratios to a Gaussian-equivalent significance.
\end{review}

The revised manuscript stops treating $\sqrt{\Delta\chi^2}$ as a detection significance for the quasar subsets and for the forecast.
Where a submitted ratio is still printed, the table heading says ``quoted ratio'' and the surrounding text says the row is not a measurement.
$R^2$ and $\chi^2$ are defined separately in Section~3.1, because the notebook's $R^2$ is unweighted and its $\chi^2$ is weighted.
AIC and BIC are not recomputed.

\revtodo{The common-likelihood AIC/BIC calculation is the same to-do as R2.7. Any future $\sigma$ equivalent needs the degrees of freedom of that likelihood written next to it.}

Change: Sections~3.1, 4.2, and~4.6.

\subsection*{R2.21}
\begin{review}
Internal numerical inconsistencies, including $H_\infty$ near $104.5$ versus high-redshift values near $71$--$72$, and the leverage and significance figures.
\end{review}

The $H_\infty$ versus $71$--$72$ clash is answered under the $10^{5}$ normalization in R2.4: they are the intercept and the finite-$t$ evaluations of one curve, on the quoted inputs.
The $\Gamma$ clash of R1.1 is not resolved.
Quoted leverage differences ($0.84$ versus $0.10$, ratio $8.4$) are retained only as the dynamic range of $\xi(t)$.
Quoted significances for the quasar subsets and the forecast are not adopted as results.
One further small inconsistency is stated: a figure caption in the submitted text printed $R^2 = 0.907$ where the abstract printed $0.91$, and the notebook prints $H_\infty = 104.469\ldots$, which rounds to $104.5$.

\revtodo{After the refits in R1.1, R2.5, and R2.7, search the manuscript again for any remaining pair of incompatible numbers. Do not silently average $6.42$ and $5.56$, or $0.907$ and $0.91$.}

Change: Sections~2.1, 4.1, and~4.2.

\subsection*{R2.22}
\begin{review}
Figures need clearer labels, units, uncertainty bands, and a separation of measurements from theory.
\end{review}

Captions of the retained figures were rewritten so that each caption states what is a quoted fit, what is an assigned anchor, and what is not a detection.
The forecast caption states that the mock has no survey covariance.
The leverage caption states that the curve is not a goodness-of-fit test.
The plotted files themselves are the submitted PDFs.
Font sizes, line widths, and axis units inside those PDFs were not regenerated.

\revtodo{Regenerate the retained figures from the checked notebooks after the refits, with readable type, units on every axis, and a legend that separates data from the curve. Do not redraw them by hand with new data values.}

Change: captions in Section~4 and in the Supplementary Material. Graphics regeneration is still a to-do.

\subsection*{R2.23}
\begin{review}
Duplicate references and bibliographic consistency.
\end{review}

The Planck 2018 entries I, V, VI, IX, and X remain distinct papers and were not deleted.
No exact duplicate bibliographic entry was removed.
The two DESI entries are different papers and are now distinguished by data release (R2.24).

\revtodo{A line-by-line check of every remaining DOI, author list, and page range is still open. This pass updated only the two DESI entries that the report singled out.}

Change: \texttt{paper/references.bib}, DESI entries. Full audit still a to-do.

\subsection*{R2.24}
\begin{review}
Reference accuracy, including DESI 2024 preprint versus DESI DR2, and missing DOIs.
\end{review}

\texttt{DESI2024} is now the Year-1 BAO paper, published as JCAP~02(2025)~021, doi \texttt{10.1088/1475-7516/2025/02/021}, with the arXiv identifier 2404.03002 retained.
\texttt{DESI2025DR2} is the DR2 BAO paper, Physical Review~D \textbf{112}, 083515 (2025), doi \texttt{10.1103/tr6y-kpc6}, arXiv 2503.14738.
The late-time DESI comparison cites the DR2 paper.
The 2024 paper is cited as Year-1 background and is not used as the DR2 measurement.
The prose anchor list that cited the 2024 paper as one of the four Hubble anchors is flagged in R2.4, because the figure notebook does not contain a BAO anchor.

Change: \texttt{references.bib}; Sections~1, 3.3, and~4.3; Supplementary DESI caption.

\subsection*{R2.25}
\begin{review}
Typographical errors, notation, and the journal footer.
\end{review}

The fitting function is written $1/\ln^2(t/t_{\mathrm{P}})$ wherever the dimensionless ratio is meant.
The Hubble equation shows the $10^{5}$ normalization explicitly.
$H_0$, $H_\infty$, $\Lambda(t)$, $\alpha(t)$, $\Gamma$, and $\chi^2/\mathrm{dof}$ are used with that notation in the revised sections.
The paragraph heading that produced ``Slowness as a physical requirement..'' has been removed.
The label ``Pillar~10'' has been removed; the letter labels that remain are identified as labels from the reviewed manuscript.
The footer of \texttt{main\_elsevier.tex} is set by \verb|\journal{Next Research}|.
The relation of $H_\infty$ to $\Lambda_\infty$ is discussed as a unit inconsistency in Section~2.1 and is not used as an equation of the fit.

\revtodo{A final proofreading pass after the refits, including a check that no unscaled $1/\ln^2(t)$ remains in a regenerated figure label.}

Change: throughout the revised sections; \texttt{main\_elsevier.tex} journal line.

\subsection*{R2.26}
\begin{review}
Conclusions should reflect the strength of the evidence.
\end{review}

Section~7 is split into three paragraphs.
The numerical paragraph lists the quoted four-anchor summary, the unresolved pair of Murphy coefficients, the opposite Keck and VLT signs, and the DESI DR2 $\chi^2/\mathrm{dof}\sim 10^{3}$.
The phenomenological paragraph limits the expansion-history reading to what the assigned pairs actually are, and it limits the laboratory statement to the conditional translation of the Rosenband et al.\ bound.
The last paragraph withdraws the readings of the curve as a solution of the Hubble tension, of the Keck coefficient as a detection, of the large late-time $\chi^2$ as support, and of the mock as an $11\sigma$ forecast.

Change: Section~7; Abstract.

\section*{What was removed, and what was kept}

Removed from the main text and the Supplementary Material: Riemann-zeta motivation and citations that existed only to support it; the lattice CMB forward model; the inverse reconstruction and held-out test; acoustic freeze-out; the post-recombination lattice visualization; three-scale and uniqueness claims; detection language for the quasar subsets; the use of $\chi^2/\mathrm{dof}\sim 10^{3}$ as support.

Kept, with the limits above: the four-anchor fit and its quoted summary (Pillar~C); the exploratory full-sample / Keck / VLT table; the dynamic range of $\xi(t)$ together with an explicit falsification standard; the laboratory derivative; the previously quoted late-time $\chi^2$ values, read as goodness of fit; the previously quoted forecast table, read as a diagonal-error mock.

We hope the revised manuscript can be reconsidered after the marked refits are supplied.
This letter and the manuscript draft have not been uploaded to the journal.

Sincerely,

Liang Wang
\end{document}
